\documentclass[11pt]{article}
\usepackage{preamble}
\usepackage{subfiles} % load last: enables per-section standalone compilation

\title{Braid Presentations for Concurrent Programs}
\author{Zachary Stone}
\date{\today}

\begin{document}
\maketitle

\begin{abstract}
Ziemiański's cube chain category $\Ch(K)$ is a combinatorial model for the space of directed paths
on a precubical set $K$. Paliga and Ziemiański built a braid invariant for precubical sets
by observing a relation between $\pi_1|\Ch(Z)|$ and braids, where $Z$ is the terminal precubical set.
However, that invariant trivializes for Euclidean $K$ because $|\Ch(\cube n)|$ is
contractible \cite[Prop.~6.1]{PZ}.
This paper constructs a localization of $\Ch(K)$ at morphisms that preserve the order
of coordinates, called $\Br(K)$. And we will show that this structure is the source of the braid
structure of $\pi_1|\Ch(Z)|$. But, there is no nerve or geometric realization in this technique.
Instead we build a 2-polygraph giving an artin braid-like presentation of
$\Br(K)$. The $\{0,1,2\}$-cells of the polygraph are the dimension $\{0,1,2\}$ chains of $K$.
Intuitively, $3$-cells of K contribute braid relations, and pairs of non-intersecting $2$-cells of
$K$ contribute commuting relations. This is not just a metaphor: at $K = Z$,
we get the artin relations and generators exactly. So $\Br(Z)$ is equivalent to the
positive braid category. The presentation is natural in $K$, providing a generalization of the
existing work. As an application, we establish a link to the hyperplane arrangement literature.
The way braids are constructed via salvetti's theorem for the braid arrangement exactly matches
the polygraph we build here for the cube.
All this work is formalized in Lean \cite{Cubical}.
\end{abstract}

\section{Introduction}\label{sec:intro}

Ziemiański's cube chain category $\Ch(K)$ models the executions of a bipointed precubical set $K$
\cite[Thm.~7.6]{Z2}, and Paliga--Ziemiański's homotopy equivalence between $|\Ch(Z)|$ and
configuration spaces of points in the plane makes it a braid invariant of any $K$, by naturality.
It trivializes on Euclidean complexes \cite[Prop.~6.1]{PZ} for a one-line reason: it factors
through the contractible $|\Ch(\cube n)|$.

The braiding is already present one level down, in $\Ch(K)$ itself, with no nerve. An object of
$\Ch(K)$ decomposes an execution into blocks of simultaneous events; a morphism is a
\emph{refinement}, a relabelling of events that coarsens the blocking. Some refinements only glue
adjacent blocks and leave every event where it was; the rest genuinely reorder events across a
block boundary. Call the first kind \emph{order preserving}, write $W$ for the class they form,
and set $\Br(K) := \Ch(K)\op[(W\op)\inv]$. Nothing is quotiented: inverting $W$ leaves a category
whose arrows remember which events crossed which, and how often. It is a braiding, not a
permutation --- two independent events that swap twice are not the identity.

We present $\Br(K)$ by generators and relations, for every $K$, naturally in $K$, and the
generators and relations are not invented: they are the low-dimensional chains of $K$. A chain of
degree $0$ is a \emph{run}, an execution with a total order on its events, and a chain $e$ of any
degree spans two runs, the run $\lorun e$ it is glued from by an order-preserving refinement and
the run $\hirun e$ that reverses every block of $e$. So chains are already arrow-shaped.

Let the \textbf{paper polygraph} $\Pgraph(K)$ (\Cref{def:poly}) have $0$-cells the runs of $K$,
$1$-cells $X \to Y$ the chains $e$ of degree $1$ with $\lorun e = X$ and $\hirun e = Y$, and
$2$-cells $X \to Y$ the chains of degree $2$ likewise, a $2$-cell $e$ being attached along the two
maximal chains of the poset of runs over $e$ (\Cref{prop:polygon}).

\begin{theorem}[\Cref{thm:A}, \Cref{thm:B}, \Cref{thm:C}]\label{thm:ABC}
  For every bipointed precubical set $K$, with no hypothesis on $K$, the polygraph $\Pgraph(K)$
  presents $\Br(K)$. The assignment $K \mapsto \Pgraph(K)$ is a functor
  $\BPSet \to \Polygraph$ and the presentation is natural in $K$, on the nose. At the terminal
  object, $\Pgraph(Z)$ is Artin's polygraph cell for cell and word for word, so
  $\Br(Z) \simeq \BraidPosOp$.
\end{theorem}

No braid theory goes in. Artin's dichotomy is the dichotomy between the two kinds of degree-two
chain: one block of dimension $3$ spans a hexagon of runs and gives
$\sigma_i\sigma_{i+1}\sigma_i = \sigma_{i+1}\sigma_i\sigma_{i+1}$, two blocks of dimension $2$ span
a square and give $\sigma_i\sigma_k = \sigma_k\sigma_i$ (\Cref{prop:polygon}). So $3$-cells of $K$
contribute braid relations and disjoint pairs of $2$-cells contribute commutations. The nearest
existing statement is Gaussent--Guiraud--Malbos's coherent presentation of an Artin monoid
\cite{GGM}, from Garside theory; here the same $2$-cells come from the geometry of $K$, and the
Artin case is the value at a point.

\paragraph{Formalisation.}
Everything below is formalised in Lean~4 on mathlib \cite{Cubical}, sorry-free; $\Br(K)$, $W$ and
$\Pgraph(K)$ are \texttt{ChainCat.Br}, \texttt{ChainCat.W} and \texttt{Paper.poly}, and
\Cref{thm:A}, \Cref{thm:B}, \Cref{thm:C} are \texttt{Paper.paperPresents},
\texttt{Paper.paperSquare} and \texttt{Paper.paperArtinIso}. The file
\texttt{CubeChains/Paper.lean} is this paper read back into the tree, each statement below
appearing there as an \texttt{example} discharged by one citation.

\paragraph{Notation.}
$\Cube$ is the box category (objects the cubes $\cube n$, morphisms the face embeddings),
$\pSet = \widehat{\Cube}$ the precubical sets, and $\BPSet$ the bipointed ones, with $Z$ final and
$\cube n$ bipointed by $(0,\dots,0)$ and $(1,\dots,1)$. We set $[n] := \{1,\dots,n\}$, write
$\cube n$ both for the cube and for the presheaf it represents, and identify
$K_n = \Hom_{\pSet}(\cube n, K)$ by Yoneda. Named ambient categories are bold ($\Cube$,
$\BraidPos$, $\Polygraph$), constructions applied to an object upright ($\Ch(K)$, $\Sq(d)$),
posets and polygraphs calligraphic ($\RunOver(d)$, $\Pgraph(K)$). Composition is classical, so $gf$
means ``first $f$, then $g$'', while a word in a polygraph is a path read left to right, so
$\sigma_i\sigma_k$ is the arrow $\sigma_k \circ \sigma_i$.

Two $\op$s are bookkeeping. Following \cite{PZ,Z2} a morphism of $\Ch(K)$ runs from a \emph{finer}
chain to a coarser one, against the direction in which an execution is built up, so we localize the
opposite category and a word of cells then composes in the reading direction; and the same
convention makes the endomorphism monoid of a loop polygraph the opposite of the monoid its words
spell, whence $\BraidPosOp$ rather than $\BraidPos$, which word reversal identifies with it.

\section{Background}\label{sec:background}

\paragraph{Cube chains \cite{PZ,Z2}.}
The \textbf{wedge} $P \vee Q$ of $P,Q \in \BPSet$ is the pushout in $\pSet$ gluing $\fin_P$ to
$\init_Q$, bipointed by $\init_P$ and $\fin_Q$; $(\BPSet,\vee,\cube 0)$ is monoidal. A
\textbf{shape} is a list $a = [a_1,\dots,a_r]$ of positive naturals, with \textbf{serial wedge}
$\bigvee a := \cube{a_1} \vee \cdots \vee \cube{a_r}$, \textbf{beads} the $\cube{a_i}$, and
$|a| := \sum_i a_i$, $\deg a := \sum_i (a_i - 1) = |a| - r$; write $\Wdg \subseteq \BPSet$ for the
full subcategory of serial wedges. Then $\Ch(K)$ is the full subcategory of $\BPSet/K$ spanned by
the serial wedges: an object is a bipointed $\bigvee a \to K$, a morphism (a \textbf{refinement})
is a map of wedges commuting over $K$, running from a finer chain to a coarser one. We write
$\deg d$, $|d|$ for the degree and event count of $d$'s shape and $\Ch(u) : \Ch(K) \to \Ch(L)$ for
the pushforward along $u : K \to L$, which postcomposes the classifying map and changes no shape. A
chain of degree $0$ is a \textbf{run}; $\Run(K) \subseteq \Ch(K)$ is the full subcategory of runs,
and it is discrete, since a refinement losing no bead is an identity.

We use three facts from \cite[\S2]{PZ}. Serial wedges are non-self-linked and admit an altitude
function $\alt : K_0 \to \Z$ with $\alt(\fin e) = \alt(\init e)+1$ on edges, the final vertex of
$\bigvee a$ having altitude $|a|$; a bipointed map between connected such objects preserves
altitude \cite[Prop.~2.7]{PZ}; and $\Ch$ is lax monoidal from $(\BPSet,\vee)$ to $(\Cat,\times)$ by
concatenation, with $\Ch(P) \times \Ch(Q) \to \Ch(P \vee Q)$ an isomorphism when $P \vee Q$ is
non-self-linked and admits altitude \cite[Lem.~2.11]{PZ}.

\paragraph{The events of a wedge.}
Let $\Coord : \Cube \to \Set$ send $\cube n$ to $[n]$ and a face embedding $g$ to the monotone
injection $g_*$ naming the directions $g$ leaves free, and let $E(\Coord)$ be its left Kan
extension along Yoneda. It is cocontinuous and $\Coord(\cube 0) = \emptyset$, so the pushout
defining the wedge becomes a coproduct:
\[
  \Ev(a) \;:=\; E(\Coord)\Bigl(\bigvee a\Bigr) \;\cong\; \coprod_i [a_i],
\]
the \textbf{events} of $a$. Write $\blk(p)$ for the bead index of an event, $f_* := E(\Coord)(f)$,
and $\str_a : \Ev(a) \cong [\,|a|\,]$ for the order isomorphism from the lexicographic order (bead
first, then direction); $\str_a(p)$ is the \textbf{strand} of $p$, its position in the canonical
execution.

\begin{proposition}\label{prop:coords}
  Let $f : \bigvee a \to \bigvee b$ in $\BPSet$. Then $f$ is determined by $f_*$, which is a
  bijection, and, writing $\blk(f)$ for the induced map on bead indices,
  \begin{equation}\label{eq:blk-coord}
    f_*(i,d) = \bigl(\blk(f)(i),\ (f_i)_*(d)\bigr), \qquad \blk \circ f_* = \blk(f) \circ \blk,
  \end{equation}
  with $\blk(f)$ monotone and each $f_i : \cube{a_i} \to \cube{b_{\blk(f)(i)}}$ a face embedding.
\end{proposition}
\begin{proof}
  By the pushout, $f$ is determined by its restrictions to the beads, and each lands in one bead of
  the target since $a_i > 0$: that gives $\blk(f)$, \eqref{eq:blk-coord} and the last clause, and a
  face embedding is determined by the directions it leaves free. Injectivity of $f_*$ and
  monotonicity of $\blk(f)$ follow by induction on the length of $b$: by \cite[Lem.~2.11]{PZ} a map
  $\bigvee a \to \cube m \vee \bigvee b'$ splits as $f_1 \vee f_2$ over a splitting
  $a = a_l \ast a_r$, so $f_* = (f_1)_* \sqcup (f_2)_*$ and
  $\blk(f) = \blk(f_1) \sqcup (1+\blk(f_2))$ with $\blk(f_1)$ constant at $1$, and injectivity of
  $(f_1)_* : \Ev(a_l) \to [m]$ is a second induction, the images at each glued vertex separated by
  the coordinates already $1$ there. Surjectivity is a count: $f$ preserves altitude
  \cite[Prop.~2.7]{PZ}, so $|a| = |b|$.
\end{proof}

\paragraph{Polygraphs, localization, and Newman's lemma.}
A \textbf{2-polygraph} $\Pgraph$ is a set $V$ of $0$-cells, families $\Pgraph_1(x,y)$,
$\Pgraph_2(x,y)$ of $1$- and $2$-cells, and two \emph{words}
$\src,\tgt : \Pgraph_2(x,y) \to \{\text{paths } x \to y\}$ in the free category on the $1$-cells,
read left to right; morphisms are maps in each dimension commuting with $\src,\tgt$ on the nose,
and $\Polygraph$ is the resulting category. The category $\pres{\Pgraph}$ \textbf{presented by}
$\Pgraph$ is the free category on the $1$-cells modulo the congruence generated by
$\src(\alpha) = \tgt(\alpha)$, with quotient functor $q$; $\Pgraph$ \textbf{presents} $C$ when
given an equivalence $\pres{\Pgraph} \to C$. We use its universal property
\cite[Ch.~2]{Polygraphs} and that of the localization \cite[Ch.~I]{GZ} in the following form.

\begin{proposition}\label{prop:present-univ}
  A functor $\pres{\Pgraph} \to C$ is the same as a map on $0$-cells together with an arrow of $C$
  per $1$-cell sending the two words of each $2$-cell to the same arrow; two such functors agreeing
  on $0$- and $1$-cells are equal; and a morphism of polygraphs bijective in every dimension is an
  isomorphism. For a class $W$ in $C$ there is $Q : C \to C[W\inv]$ inverting $W$, through which
  every functor inverting $W$ factors uniquely, and two functors out of $C[W\inv]$ agreeing after
  $Q$ are equal.
\end{proposition}

For $m : a \to e$ in $W$ and $f : b \to e$ arbitrary we write
$\mathrm{frac}(m,f) := Q(f\op) \circ Q(m\op)\inv : Q(a) \to Q(b)$ in $C\op[(W\op)\inv]$, ``back
along the order-preserving leg, forward along the other''. Then $\mathrm{frac}(m,m) = \id$ and two
such glue: for $l : e_0 \to e$, $m_0 : b \to e_0$ in $W$ and $f_0 : c \to e_0$,
\begin{equation}\label{eq:frac-glue}
  \mathrm{frac}(m_0,f_0) \circ \mathrm{frac}(m,\ l \circ m_0) \;=\; \mathrm{frac}(m,\ l \circ f_0).
\end{equation}

Write $\mathrm{Hasse}(P)$ for the Hasse quiver of a poset $P$, and call a map sending each element
to an object of $C$ and each cover to an arrow a \textbf{reading}; it extends to paths.

\begin{theorem}[coherent Newman; \cite{Huet}, {\cite[Ch.~4]{Polygraphs}}]\label{thm:newman}
  Let $P$ be finite and $\varphi$ a reading of $P$ in $C$ such that for all covers
  $x \lessdot y_1$, $x \lessdot y_2$ with $y_1 \neq y_2$ and every common upper bound $z$ of
  $y_1,y_2$, there are $j \le z$ and paths $p_i : y_i \to j$ with
  $\varphi(x \lessdot y_1)\varphi(p_1) = \varphi(x \lessdot y_2)\varphi(p_2)$. Then $\varphi$ sends
  parallel paths to equal arrows, so it extends uniquely to a functor $P \to C$, and a functor out
  of $P$ is pinned by its values on covers.
\end{theorem}

\begin{lemma}\label{lem:polygon}
  If a finite poset $P$ is graded by $r$ with $r(y) = r(x)+1$ on covers and every element has
  exactly two Hasse neighbours, then a directed path is determined by its first edge; and if $P$
  has a least and a greatest element then $P$ is a cycle, its two maximal chains have a common
  length $\ell$, and $|P| = 2\ell$. (Undirected, the Hasse graph is $2$-regular hence a union of
  cycles; $r$ increases along a path, so it leaves each vertex by the edge it did not arrive on;
  both edges at $\bot$ point up and both at $\top$ point down, so the cycle through $\bot$ meets
  $\top$ and splits into two arcs of equal length, and it is all of $P$.)
\end{lemma}

\section{Squares, runs, and the braiding}\label{sec:squares}

\subsection{Squares, triangles, and biclosed sets}

\begin{definition}\label{def:squares}
  Let $\Sq, \Tri : \Cube \to \Set$ send $\cube n$ to the monotone injections
  $[2] \hookrightarrow [n]$, respectively $[3] \hookrightarrow [n]$. Both vanish at $\cube 0$, so
  their left Kan extensions along Yoneda give, as for $\Ev$,
  \[
    \Sq(a) \;\cong\; \coprod_i \binom{[a_i]}{2}, \qquad \Tri(a) \;\cong\; \coprod_i \binom{[a_i]}{3},
    \qquad\text{so}\qquad |\Sq(a)| = \sum_i \binom{a_i}{2}.
  \]
  Concretely a \textbf{square} $s$ is a pair of events of one bead, written
  $\mathrm{lo}(s), \mathrm{hi}(s)$ with $\str(\mathrm{lo}\,s) < \str(\mathrm{hi}\,s)$, and a
  \textbf{triangle} $t$ is a triple of events of one bead, with edges
  $t_{01}, t_{02}, t_{12} \in \Sq(a)$. We write $\Sq(d)$ for the squares of a chain's shape and
  $q_* : \Sq(o) \to \Sq(d)$ for the map induced by a refinement, injective by
  \Cref{prop:coords}. A set $R \subseteq \Sq(a)$ is \textbf{biclosed} when for every triangle $t$,
  \[
    t_{01},t_{12} \in R \implies t_{02} \in R,
    \qquad\text{and}\qquad
    t_{02} \in R \implies t_{01} \in R \text{ or } t_{12} \in R;
  \]
  equivalently, both $R$ and $R^c$ are closed under the first rule.
\end{definition}

\begin{lemma}\label{lem:sq-range}
  For $q : \bigvee a \to \bigvee b$, a square $s$ of $\bigvee b$ lies in $\im q_*$ iff
  $q_*\inv(\mathrm{lo}\,s)$ and $q_*\inv(\mathrm{hi}\,s)$ share a bead of $\bigvee a$; and for
  events $u,v$ of $\bigvee b$ in distinct beads,
  $\blk(q_*\inv u) < \blk(q_*\inv v) \iff \blk(u) < \blk(v)$.
\end{lemma}
\begin{proof}
  If $s = q_*t$ the events pull back to $t$'s, which share a bead; conversely if they share a bead
  $i$ then, $(q_i)_*$ being a monotone injection, they are the events of a square of that bead. The
  second claim is monotonicity of $\blk(q)$ (\Cref{prop:coords}) in both directions, the beads
  being distinct.
\end{proof}

\subsection{Runs of a wedge}

A run over the wedge $\bigvee a$ is a chain $x : \bigvee[1^N] \to \bigvee a$ of degree $0$, where
$N = |a|$; by \Cref{prop:coords} it is determined by the bijection
$x_* : [N] = \Ev([1^N]) \to \Ev(a)$. Write $\step_x := x_*\inv$ for its \textbf{firing order} and
\[
  \rev(x) \;:=\; \{\, s \in \Sq(a) : \step_x(\mathrm{hi}\,s) < \step_x(\mathrm{lo}\,s) \,\}
\]
for the squares it \textbf{reverses}.

\begin{lemma}\label{lem:firing-blocks}
  A run fires the beads in blocks: $\blk(p) < \blk(q)$ implies $\step_x(p) < \step_x(q)$. So a run
  is exactly a choice of linear order on $[a_i]$ for each bead, and $\blk$ is a monotone surjection
  from the firing order onto the beads, with interval fibres.
\end{lemma}
\begin{proof}
  The beads of $[1^N]$ are singletons, so $\blk$ there is the identity and
  $\blk(x) = \blk \circ x_*$, monotone by \Cref{prop:coords}. A monotone surjection from a chain
  has interval fibres, here the beads; the only freedom left is the order inside each.
\end{proof}

\begin{theorem}[a run \emph{is} a biclosed set]\label{thm:structure}
  $x \mapsto \rev(x)$ is a bijection from the runs of $\bigvee a$ onto the biclosed subsets of
  $\Sq(a)$. Reversing every bead's order complements $\rev$; in particular biclosed sets are closed
  under complement.
\end{theorem}
\begin{proof}
  A triangle lies in one bead, so by \Cref{lem:firing-blocks} the statement is bead-local; fix a
  bead $[m]$ with its cube order $<$. A linear order $\prec$ on $[m]$ has inversion set
  $\rev(\prec) = \{(p,q) : p<q,\ q \prec p\}$, from which it is recovered by: $p \prec q$ iff
  ($p<q$ and $(p,q) \notin R$) or ($q<p$ and $(q,p) \in R$). So $\rev$ is injective, and its image
  is the set of $R$ for which that recipe is transitive --- totality being automatic. Transitivity
  can only fail on a triple $p<q<r$, where the cases to check are exactly the two biclosedness
  clauses at that triangle. Complementation reverses each $\prec$.
\end{proof}

For one bead of dimension $m$ this is the classical description of inversion sets in $S_m$, and
inclusion of reversed sets is the right weak order \cite[\S3.1]{BB}; the order theory below is that
statement, bead by bead.

\begin{lemma}\label{lem:lower-biclosed}
  Order $\Sq(a)$ by $\str(\mathrm{lo}\,-)$, then by $\str(\mathrm{hi}\,-)$; call this the
  \textbf{standard order}. Its lower sets and its upper sets are all biclosed.
\end{lemma}
\begin{proof}
  On a triangle with events $p<q<r$ the edges are $t_{01}=(p,q)$, $t_{02}=(p,r)$, $t_{12}=(q,r)$,
  in that standard order. For a lower set $R$: $t_{12} \in R$ forces $t_{02} \in R$ and
  $t_{02} \in R$ forces $t_{01} \in R$, which are the two clauses. For an upper set: $t_{01} \in R$
  forces $t_{02} \in R$ and $t_{02} \in R$ forces $t_{12} \in R$.
\end{proof}

\subsection{Crossings, order-preserving refinements, and \texorpdfstring{$\Br(K)$}{Br(K)}}

\begin{definition}\label{def:crossings}
  A map of serial wedges $\varphi : \bigvee a \to \bigvee b$ \textbf{crosses}
  \[
    \Cross(\varphi) \;:=\;
      \bigl\{\, s \in \Sq(b) : \blk\bigl(\varphi_*\inv(\mathrm{hi}\,s)\bigr)
      < \blk\bigl(\varphi_*\inv(\mathrm{lo}\,s)\bigr) \,\bigr\},
  \]
  the squares of the target whose upper event it brings from a strictly earlier bead than their
  lower. For a refinement we write $\Cross(f)$ for the crossings of its wedge map.
\end{definition}

\begin{lemma}[the cocycle]\label{lem:cocycle}
  $\Cross(\psi\varphi) = \psi_*(\Cross\varphi) \sqcup \Cross(\psi)$, disjointly, for
  $\varphi : \bigvee a \to \bigvee b$ and $\psi : \bigvee b \to \bigvee c$; and
  $\Cross(\id) = \emptyset$.
\end{lemma}
\begin{proof}
  Take $s \in \Sq(c)$. If $s = \psi_*(t)$ then $s \notin \Cross(\psi)$, since $t$'s events share a
  bead (\Cref{lem:sq-range}) --- which also gives disjointness --- and
  $(\psi\varphi)_*\inv$ agrees with $\varphi_*\inv$ on $t$'s events, so
  $s \in \Cross(\psi\varphi)$ iff $t \in \Cross(\varphi)$. If $s \notin \im\psi_*$ then
  $s \notin \psi_*(\Cross\varphi)$ and, the events $\psi_*\inv(\mathrm{lo}\,s)$,
  $\psi_*\inv(\mathrm{hi}\,s)$ lying in distinct beads, the second part of \Cref{lem:sq-range}
  applied to $\varphi$ gives $s \in \Cross(\psi\varphi)$ iff $s \in \Cross(\psi)$.
\end{proof}

\begin{definition}\label{def:W}
  A refinement $f$ is \textbf{order preserving} when $\Cross(f) = \emptyset$; write $W = W_K$ for
  the class of these, and
  \[
    \Br(K) \;:=\; \Ch(K)\op\bigl[(W\op)\inv\bigr], \qquad [c] := Q(c).
  \]
  Since $W$ is a condition on the wedge map alone, $\Ch(u)$ carries $W_K$ into $W_L$ and descends
  to $\Br(u) : \Br(K) \to \Br(L)$, strictly functorially in $u$.
\end{definition}

By \Cref{prop:W} below, $f \in W$ means $f$ only glues adjacent beads; the elementary ones are the
comparisons $\bigvee[\dots,p,q,\dots] \to \bigvee[\dots,p+q,\dots]$, and $W$ is the multiplicative
class they generate.

\begin{proposition}\label{prop:W}
  For a refinement $f : c \to d$ the following agree: $f \in W$; $f_*$ is monotone for the strand
  orders; $\str_d \circ f_* = \str_c$, i.e.\ every event keeps its position. Consequently $W$
  contains the identities, $fg \in W$ iff $f \in W$ and $g \in W$, and two order-preserving
  refinements with the same source and target are equal.
\end{proposition}
\begin{proof}
  A monotone bijection of finite chains of equal size is the unique isomorphism, which gives the
  second and third conditions from each other. For the first two, take $u,v \in \Ev(c)$ with
  $\str(u) < \str(v)$. If $\blk(u) = \blk(v)$ then $f_*u, f_*v$ lie in one bead in the same order,
  $(f_i)_*$ being a monotone injection. If $\blk(u) < \blk(v)$ then $\blk(f_*u) \le \blk(f_*v)$, so
  $\str(f_*u) > \str(f_*v)$ can happen only when both land in one bead, and then
  $s := (f_*v, f_*u)$ is a square with $\blk(f_*\inv \mathrm{hi}\,s) < \blk(f_*\inv \mathrm{lo}\,s)$,
  that is $s \in \Cross(f)$; conversely every crossing arises so. Closure and two-out-of-three are
  \Cref{lem:cocycle} with $\psi_*$ injective, and the last claim holds because the third condition
  pins $f_* = \str_d\inv \str_c$, which pins $f$ (\Cref{prop:coords}).
\end{proof}


\section{The paper polygraph}\label{sec:poly}

Fix $K$, and let $d$ be a chain of $K$ of shape $a$, with $N = |d|$.

\subsection{Runs over a chain}

\begin{definition}\label{def:runover}
  A \textbf{run over $d$} is a run $X$ of $K$ with a refinement $X \to d$; write $\RunOver(d)$ for
  the set of them, $\mathrm{run}(x)$ and $\mathrm{ref}(x)$ for the two components, and
  $q_\sharp x := q \circ \mathrm{ref}(x)$ for the \textbf{push} along a refinement $q : o \to d$.
\end{definition}

\begin{lemma}\label{lem:runover}
  Composing with $d$'s classifying map identifies $\RunOver(d)$ with the runs of $\bigvee a$, and
  under that identification $\rev(x) = \Cross(\mathrm{ref}(x))$. So, ordering $\RunOver(d)$ by
  $x \le y :\!\iff \rev(x) \subseteq \rev(y)$, the map $\rev$ is an isomorphism of posets onto the
  biclosed subsets of $\Sq(d)$, with least element $\bot_d$ ($\rev = \emptyset$), greatest element
  $\top_d$ ($\rev = \Sq(d)$), and complementation as an anti-automorphism. Write
  $\lorun d, \hirun d \in \Run(K)$ for the runs underlying $\bot_d$ and $\top_d$.
\end{lemma}
\begin{proof}
  A refinement $X \to d$ over $K$ is a wedge map $\bigvee[1^N] \to \bigvee a$ whose composite with
  $d$'s classifying map is $X$'s, so the latter is forced and the data is just the wedge map. Every
  bead of the source is then an edge, so $\blk$ on the source is $\str$ and
  $\blk \circ \mathrm{ref}(x)_*\inv = \step_x$, which makes \Cref{def:crossings} and
  the display above the same condition. The rest is \Cref{thm:structure}.
\end{proof}

\begin{proposition}[the push formula]\label{prop:push}
  $\rev(q_\sharp x) = q_*(\rev x) \sqcup \Cross(q)$. Hence $q_*\inv\rev(q_\sharp x) = \rev(x)$, so
  $q_\sharp$ is an order embedding; and $q_\sharp\bot_o = \bot_d$ iff $q \in W$. Moreover, for
  $x \in \RunOver(d)$ the conditions $\mathrm{ref}(x) \in W$, $\rev(x) = \emptyset$ and
  $x = \bot_d$ agree, so $\lorun X = X$ for a run $X$, and $\lorun d = \lorun c$ whenever
  $c \to d$ lies in $W$.
\end{proposition}
\begin{proof}
  The formula is \Cref{lem:cocycle} read through \Cref{lem:runover}. Since $q_*$ is injective with
  image missing $\Cross(q)$ (\Cref{lem:sq-range}), it follows that
  $q_*\inv\rev(q_\sharp x) = \rev(x)$ and that $q_\sharp x \le q_\sharp y$ iff $x \le y$; and
  $\rev(q_\sharp\bot_o) = \Cross(q)$. The last sentence is \Cref{lem:runover} and the formula
  applied to $c \to d$.
\end{proof}

\subsection{Steps, covers, and cells}

\begin{lemma}\label{lem:steps}
  Call $s \in \Sq(d)$ a \textbf{step} of $x \in \RunOver(d)$ when $\rev(x) \symdiff \{s\}$ is
  biclosed, write $\Steps(x)$ for the set of them and $x^s$ for the \textbf{flip}, the run
  reversing $\rev(x) \symdiff \{s\}$. Then $s$ is a step iff its two events are consecutive in the
  firing order of $x$, and $|\Steps(x)| = \deg d$.
\end{lemma}
\begin{proof}
  A bead's events are consecutive in the firing order (\Cref{lem:firing-blocks}), so ``consecutive
  in the bead'' and ``consecutive overall'' agree. If the two events are consecutive, toggling $s$
  swaps an adjacent pair of a linear order, so the result is biclosed by \Cref{thm:structure}; if
  some event $w$ of the bead fires strictly between them, toggling swaps them but not their order
  against $w$, so $w$ would lie both before and after the pair. Thus the steps are the covers of
  the firing order lying inside a bead. That order is a chain on $N$ elements and $\blk$ is a
  monotone surjection onto the $r$ beads (\Cref{lem:firing-blocks}); such a map collapses exactly
  $N - r$ covers, so there are $N - r = \deg d$ of them.
\end{proof}

\begin{lemma}[the cover lemma]\label{lem:cover}
  If $x < z$ then some $s \in \Steps(x) \setminus \rev(x)$ lies in $\rev(z)$. Consequently
  $x \lessdot y$ iff $\rev(y) = \rev(x) \cup \{s\}$ for such an $s$, in which case $y = x^s$; so
  $\RunOver(d)$ is graded by $|\rev|$, from $0$ at $\bot_d$ to $|\Sq(d)|$ at $\top_d$.
\end{lemma}
\begin{proof}
  Write $\prec_x, \prec_z$ for the firing orders. If $z$ agreed with $x$ on every $x$-covering
  pair then $\prec_x \subseteq \prec_z$ by transitivity, so $z = x$. So some $x$-cover
  $p \lessdot_x q$ has $q \prec_z p$; its events share a bead, since otherwise $\blk(p) < \blk(q)$
  and every run fires $p$ first (\Cref{lem:firing-blocks}). So the pair is a step $s$ of $x$
  (\Cref{lem:steps}); if $\mathrm{lo}\,s = p$ then $s \in \rev(z) \setminus \rev(x)$ as wanted,
  and if $\mathrm{lo}\,s = q$ then $s \in \rev(x) \subseteq \rev(z)$ says $p \prec_z q$, a
  contradiction. For the covering criterion, $x \le x^s \le y$ forces $y = x^s$ in one direction,
  and in the other two biclosed sets differing by one element are a covering pair under
  \Cref{lem:runover}.
\end{proof}

\begin{proposition}[a factorisation is a set of steps]\label{prop:factorisations}
  Sending a factorisation $\mathrm{run}(x) \to m \to d$ of $\mathrm{ref}(x)$ to the set of steps of
  $x$ whose events share a bead of $m$ is a bijection onto the subsets of $\Steps(x)$, with
  $\deg m = |S|$. Write $m_S$ for the middle and $\iota_S, \pi_S$ for the legs. If moreover
  $S \cap \rev(x) = \emptyset$ --- we then say $S$ \textbf{rises} from $x$ --- then $\iota_S \in W$,
  so $\lorun{m_S} = \mathrm{run}(x)$.
\end{proposition}
\begin{proof}
  A middle $m$ makes the $N$ events into its beads, so it cuts the firing order of $x$ into
  intervals (\Cref{lem:firing-blocks} for $x$ over $m$), and $\pi$ requires each interval to sit
  inside one bead of $d$. Conversely such an interval partition arises from the chain whose beads
  are its blocks in order; and the factorisation is pinned by that shape, because a wedge map is
  determined by its action on events and injective there (\Cref{prop:coords}), hence monic, while
  the first leg is forced by which events fall in which block. An interval partition of a chain is
  the set of covers it does not cut, and those lie inside a bead, so they are steps
  (\Cref{lem:steps}); and $\deg m = N - \#\text{blocks} = \#\text{uncut covers} = |S|$.

  If $S$ rises, then by \Cref{lem:runover} $\Cross(\iota_S)$ is the set of squares of $m_S$ that
  $x$ reverses, that is the pairs inside one block that $x$ reverses. The covers inside a block all
  lie in $S$, so none is reversed, so the firing order agrees with the strand order on consecutive
  pairs of the block and hence, by transitivity, on all of them. So $\Cross(\iota_S) = \emptyset$.
\end{proof}

\begin{definition}\label{def:cell}
  An \textbf{$n$-cell} $X \to Y$ is a chain $e$ with $\deg e = n$, $\lorun e = X$ and
  $\hirun e = Y$; write $\Cell_n(X,Y)$. So the $n$-cells are just the chains of degree $n$, sorted
  by the two runs each spans.
\end{definition}

\begin{lemma}[the closure lemma]\label{lem:closure}
  Let $S$ rise from $x$. Then
  $\rev\bigl((\pi_S)_\sharp\top_{m_S}\bigr) = \rev(x) \cup (\pi_S)_*\bigl(\Sq(m_S)\bigr)$, and any
  $z \ge x$ with $S \subseteq \rev(z)$ has $(\pi_S)_*(\Sq(m_S)) \subseteq \rev(z)$. So
  $(\pi_S)_\sharp\top_{m_S}$ is the least element above $x$ reversing all of $S$.
\end{lemma}
\begin{proof}
  The push formula gives
  $\rev((\pi_S)_\sharp\top_{m_S}) = (\pi_S)_*\Sq(m_S) \cup \Cross(\pi_S)$, and the cocycle applied
  to $\pi_S \iota_S = \mathrm{ref}(x)$ with $\Cross(\iota_S) = \emptyset$ gives
  $\Cross(\pi_S) = \rev(x)$. For the second claim, list a block as $p_1 < \cdots < p_k$ in strand
  order --- which is its firing order, as in the previous proof. Each $(p_i,p_{i+1})$ lies in $S$,
  so $z$ reverses it; since $\rev(z)$ is biclosed, the first clause at the triangles
  $(p_i,p_{i+1},p_j)$ gives by induction that $z$ reverses every pair of the block.
\end{proof}

\begin{corollary}[covers are $1$-cells]\label{cor:covercell}
  For a cover $x \lessdot y$ with step $s$, the chain $\gamma(x \lessdot y) := m_{\{s\}}$ is a
  $1$-cell $\mathrm{run}(x) \to \mathrm{run}(y)$, with
  $(\pi_{\{s\}})_\sharp\bot_{m_{\{s\}}} = x$ and $(\pi_{\{s\}})_\sharp\top_{m_{\{s\}}} = y$;
  conversely every $1$-cell $\alpha$ is $\gamma$ of the unique cover
  $\bot_\alpha \lessdot \top_\alpha$ of $\RunOver(\alpha)$. So $\gamma$, with
  $x \mapsto \mathrm{run}(x)$ on vertices, is a map of quivers
  $\mathrm{Hasse}(\RunOver(d)) \to (\Run(K), \Cell_1)$, and it is natural:
  $\gamma(q_\sharp e) = \gamma(e)$ for every refinement $q : o \to d$.
\end{corollary}
\begin{proof}
  $m_{\{s\}}$ has degree $1$ and $\lorun{m_{\{s\}}} = \mathrm{run}(x)$ by
  \Cref{prop:factorisations}; by \Cref{lem:closure} its top pushes to the least element above $x$
  reversing $s$, namely $x^s = y$. Conversely a chain $\alpha$ of degree $1$ has $|\Sq(\alpha)| = 1$,
  so $\RunOver(\alpha) = \{\bot_\alpha \lessdot \top_\alpha\}$ and the factorisation through that
  one step is $\alpha$ itself. For naturality, the step of $q_\sharp e$ is $q_*$ of the step of $e$
  (\Cref{prop:push}), and $q_*$ carries the firing order of $x$ to that of $q_\sharp x$, so the two
  factorisations cut the same partition and have the same middle.
\end{proof}

\subsection{The attachment of a \texorpdfstring{$2$}{2}-cell, and the polygraph}

\begin{lemma}\label{lem:orderpath}
  If $\prec$ is a linear order on $\Sq(d)$ whose lower sets are all biclosed, the runs $x_t$
  reversing its first $t$ squares form a maximal chain
  $\bot_d = x_0 \lessdot \cdots \lessdot x_{|\Sq(d)|} = \top_d$, and the $(t+1)$-st square is a
  step of $x_t$. By \Cref{lem:lower-biclosed} the standard order and its dual both qualify.
\end{lemma}
\begin{proof}
  Each $x_t$ exists by \Cref{lem:runover}, consecutive ones differ by one element, and
  \Cref{lem:cover} applies.
\end{proof}

\begin{proposition}[the polygon, and the two species]\label{prop:polygon}
  Let $\deg e = 2$. Then every $x \in \RunOver(e)$ has exactly two steps, hence exactly two Hasse
  neighbours; $\RunOver(e)$ is a polygon with $2|\Sq(e)|$ elements; and its only two maximal chains
  from $\bot_e$ to $\top_e$ are the sweeps of the standard order and of its dual. Write $\src(e)$
  and $\tgt(e)$ for the words of $1$-cells they read off under $\gamma$; they have length
  $|\Sq(e)|$ and differ in their first letter. Moreover $e$ has either one bead of dimension $3$,
  with $|\Sq(e)| = 3$ and $\RunOver(e)$ a hexagon, or two beads of dimension $2$, with
  $|\Sq(e)| = 2$ and $\RunOver(e)$ a square.
\end{proposition}
\begin{proof}
  $|\Steps(x)| = 2$ by \Cref{lem:steps}, and the two flips are the two neighbours
  (\Cref{lem:cover}, up or down according to whether the step is reversed). The two sweeps of
  \Cref{lem:orderpath} leave $\bot_e$ by different covers, the least and greatest squares of the
  standard order being distinct, so \Cref{lem:polygon} applies with the grading $|\rev|$. Finally
  $\deg e = 2$ forces the bead dimensions above $1$ to be $\{3\}$ or $\{2,2\}$, and
  \Cref{def:squares} counts $\binom32 = 3$ and $2\binom22 = 2$.
\end{proof}

\begin{definition}\label{def:poly}
  The \textbf{paper polygraph} $\Pgraph(K)$ has $0$-cells $\Run(K)$, which is a set since a
  refinement between chains of degree $0$ is an identity; $1$-cells $\Cell_1(X,Y)$; $2$-cells
  $\Cell_2(X,Y)$; and the words $\src, \tgt$ of \Cref{prop:polygon}. Its data is nothing but the
  chains of $K$ of degree $0$, $1$ and $2$; the only computation is the attachment, and it has
  exactly two answers.
\end{definition}

\begin{proposition}[diverging covers close]\label{prop:closes}
  Read in $\pres{\Pgraph(K)}$, the quiver map $\gamma$ satisfies the hypothesis of
  \Cref{thm:newman} on $\RunOver(d)$, so it extends uniquely to a functor
  $\Gamma_d : \RunOver(d) \to \pres{\Pgraph(K)}$, and $\Gamma_d \circ q_\sharp = \Gamma_o$ for
  every refinement $q : o \to d$.
\end{proposition}
\begin{proof}
  Let $x \lessdot y_1$, $x \lessdot y_2$ be distinct covers, with steps $s_1 \neq s_2$, and $z$ a
  common upper bound. Then $S = \{s_1,s_2\}$ rises from $x$, so $\alpha := m_S$ is a chain of
  degree $2$, that is a $2$-cell (\Cref{prop:factorisations}), and $j := (\pi_S)_\sharp\top_\alpha$
  satisfies $j \le z$ by \Cref{lem:closure}. By \Cref{prop:polygon} the two maximal chains of
  $\RunOver(\alpha)$ are $\src(\alpha)$ and $\tgt(\alpha)$ and leave $\bot_\alpha$ by its two
  covers. Pushing along $\pi_S$ carries $\bot_\alpha$ to $x$, $\top_\alpha$ to $j$, covers to
  covers (\Cref{prop:push}, \Cref{lem:cover}) and keeps each cover's $1$-cell
  (\Cref{cor:covercell}); so the two chains push to two paths $x \to j$ leaving by $y_1$ and by
  $y_2$ and reading as $\src(\alpha)$ and $\tgt(\alpha)$, equal in $\pres{\Pgraph(K)}$. For
  naturality, $q_\sharp$ is monotone, so $\Gamma_d \circ q_\sharp$ and $\Gamma_o$ are functors out
  of $\RunOver(o)$ agreeing on covers, hence equal by \Cref{thm:newman}.
\end{proof}

\section{Theorem A: the polygraph presents \texorpdfstring{$\Br(K)$}{Br(K)}}\label{sec:thmA}

Theorem A compares two universal properties: a functor out of $\pres{\Pgraph(K)}$ is a reading of
the degree-one chains identifying the two words of every degree-two chain
(\Cref{prop:present-univ}), and a functor out of $\Br(K)$ is a functor on $\Ch(K)\op$ inverting $W$
(\Cref{prop:present-univ}). We build each from the other.

\begin{definition}\label{def:E}
  A cell $\alpha \in \Cell_n(X,Y)$ carries two legs, the refinements of its extreme runs over
  itself: $\alpha_\bot := \mathrm{ref}(\bot_\alpha) : X \to \alpha$, which is order preserving, and
  $\alpha_\top := \mathrm{ref}(\top_\alpha) : Y \to \alpha$. Put $E(X) := [X]$ and
  $E(\alpha) := \mathrm{frac}(\alpha_\bot, \alpha_\top) : [X] \to [Y]$.
\end{definition}

\begin{lemma}[telescoping, and soundness]\label{lem:telescope}
  For $z \in \RunOver(d)$ put
  $\mathrm{at}(z) := \mathrm{frac}\bigl(\mathrm{ref}(\bot_d), \mathrm{ref}(z)\bigr)$. Then
  $\mathrm{at}(\bot_d) = \id$ and $\mathrm{at}(y) = E(\gamma(x \lessdot y)) \circ \mathrm{at}(x)$
  for every cover, so every path from $\bot_d$ to $z$ reads as $\mathrm{at}(z)$ under
  $E \circ \gamma$. In particular the two words of a $2$-cell $e$ both read as
  $\mathrm{at}(\top_e)$, so $E$ extends to a functor $E : \pres{\Pgraph(K)} \to \Br(K)$.
\end{lemma}
\begin{proof}
  $\mathrm{at}(\bot_d) = \mathrm{frac}(m,m) = \id$. For a cover with step $s$, put
  $\beta := \gamma(x \lessdot y)$ and $q := \pi_{\{s\}}$; by \Cref{cor:covercell},
  $\mathrm{ref}(x) = q\beta_\bot$ and $\mathrm{ref}(y) = q\beta_\top$, so \eqref{eq:frac-glue} with
  $l = q$ gives $E(\beta) \circ \mathrm{at}(x) = \mathrm{frac}(\mathrm{ref}(\bot_d), q\beta_\top)
  = \mathrm{at}(y)$. Both words of a $2$-cell are paths $\bot_e \to \top_e$
  (\Cref{prop:polygon}); now apply \Cref{prop:present-univ}.
\end{proof}

\begin{definition}\label{def:theta}
  Define $\Theta : \Ch(K)\op \to \pres{\Pgraph(K)}$ by $\Theta(d) := \lorun d$ and, for
  $u : c \to d$ in $\Ch(K)$, $\Theta(u\op) := \Gamma_d(\bot_d \le u_\sharp\bot_c)$, an arrow
  $\lorun d \to \lorun c$.
\end{definition}

\begin{lemma}\label{lem:theta}
  $\Theta$ is a functor and inverts $W\op$; indeed $\Theta(u\op)$ is an identity for $u \in W$. So
  it descends uniquely to $\Phi : \Br(K) \to \pres{\Pgraph(K)}$ with $\Phi Q = \Theta$.
\end{lemma}
\begin{proof}
  For $v : b \to c$ and $u : c \to d$ we have
  $\bot_d \le u_\sharp\bot_c \le u_\sharp v_\sharp\bot_b = (uv)_\sharp\bot_b$, using that
  $u_\sharp$ is monotone (\Cref{prop:push}); since $\Gamma_d$ is a functor, $\Theta((uv)\op)$ is
  $\Gamma_d(u_\sharp\bot_c \le u_\sharp v_\sharp\bot_b) \circ \Theta(u\op)$, and the left factor is
  $\Gamma_c(\bot_c \le v_\sharp\bot_b) = \Theta(v\op)$ by naturality (\Cref{prop:closes}). For
  $u \in W$, $u_\sharp\bot_c = \bot_d$ (\Cref{prop:push}). Now apply \Cref{prop:present-univ}.
\end{proof}

\begin{theorem}[Theorem A]\label{thm:A}
  For every $K$, with no hypothesis on $K$, the functor $E$ is an equivalence
  $\pres{\Pgraph(K)} \simeq \Br(K)$, with inverse $\Phi$. So the runs of $K$ and its chains of
  degree one and two present $\Br(K)$. It is an equivalence and no more: $\Br(K)$ is not skeletal,
  a chain and the run below it becoming isomorphic there without being equal, and the runs are a
  skeleton.
\end{theorem}
\begin{proof}
  \emph{$\Phi E = \id$.} On a $0$-cell, $\Phi(E(X)) = \Theta(X) = \lorun X = X$
  (\Cref{prop:push}). On a $1$-cell $\alpha : X \to Y$, reading a cospan back gives
  $\Phi(\mathrm{frac}(m,f)) = \Theta(f\op)\Theta(m\op)\inv$, and $\Theta(m\op)$ is an identity for
  $m \in W$, so $\Phi(E(\alpha)) = \Theta(\alpha_\top\op)$. Now $(\alpha_\top)_\sharp\bot_Y$ is
  $\alpha_\top$ itself, that is $\top_\alpha$, and $\RunOver(\alpha)$ is the two-element poset
  $\bot_\alpha \lessdot \top_\alpha$, so
  $\Theta(\alpha_\top\op) = \Gamma_\alpha(\bot_\alpha \le \top_\alpha)
   = q(\gamma(\bot_\alpha \lessdot \top_\alpha)) = q(\alpha)$ by \Cref{cor:covercell}. Conclude by
  \Cref{prop:present-univ}.

  \emph{$E\Phi \cong \id$.} By \Cref{lem:telescope},
  $E(\Theta(u\op)) = \mathrm{at}(u_\sharp\bot_c)
   = \mathrm{frac}(\mathrm{ref}(\bot_d), u \circ \mathrm{ref}(\bot_c))$, which unfolds to
  $Q(\mathrm{ref}(\bot_c)\op) \circ Q(u\op) \circ Q(\mathrm{ref}(\bot_d)\op)\inv$. So
  $E\Theta \cong Q$, with component at $c$ the isomorphism $[c] \cong [\lorun c]$ that
  $\mathrm{ref}(\bot_c) \in W$ becomes in the localization; by uniqueness of descent
  (\Cref{prop:present-univ}), $E\Phi \cong \id_{\Br(K)}$.
\end{proof}

\section{Theorem B: naturality in \texorpdfstring{$K$}{K}}\label{sec:thmB}

\begin{theorem}[Theorem B]\label{thm:B}
  $K \mapsto \Pgraph(K)$, with $\Pgraph(f) := \Ch(f)$ on cells in every dimension, is a functor
  $\BPSet \to \Polygraph$, and for every $f : K \to L$ the presentation square commutes as an
  equality of functors:
  \[
    E_L \circ \pres{\Pgraph(f)} \;=\; \Br(f) \circ E_K
    \;:\; \pres{\Pgraph(K)} \longrightarrow \Br(L).
  \]
\end{theorem}
\begin{proof}
  Write $f_\bullet := \Ch(f)$. A run over a chain $e$ is a run of $e$'s shape (\Cref{lem:runover}),
  and $f_\bullet e$ has that same shape with the classifying map forced on both sides, so
  $\RunOver(e) = \RunOver(f_\bullet e)$ as ordered sets --- and everything in \S\ref{sec:squares},
  \Cref{lem:steps} and \Cref{cor:covercell} is defined from the wedge map alone. Hence $f_\bullet$
  preserves $\lorun{\,\cdot\,}$, $\hirun{\,\cdot\,}$, cells, steps, covers and $\gamma$, and it
  carries the sweep of the standard order on $\Sq(e)$ to the sweep of the standard order on
  $\Sq(f_\bullet e) = \Sq(e)$, so it preserves $\src$ and $\tgt$ letter by letter: $\Pgraph(f)$ is
  a morphism of polygraphs, functorially in $f$ because $\Ch$ is a functor.

  Both composites send $X$ to $[f_\bullet X]$. On a $1$-cell $\alpha$, $\Br(f)$ is the descent of
  $Q_L \circ f_\bullet\op$, so it reads a cospan upstairs:
  $\Br(f)\mathrm{frac}(m,g) = \mathrm{frac}(f_\bullet m, f_\bullet g)$. Here $f_\bullet\alpha_\bot$
  is an order-preserving refinement $f_\bullet X \to f_\bullet\alpha$, hence equals
  $(f_\bullet\alpha)_\bot$ by \Cref{prop:W}, and $f_\bullet\alpha_\top = (f_\bullet\alpha)_\top$
  likewise. So both composites send $\alpha$ to
  $\mathrm{frac}((f_\bullet\alpha)_\bot, (f_\bullet\alpha)_\top)$, and \Cref{prop:present-univ}
  finishes. The square commutes strictly --- more than a presentation is entitled to --- because an
  order-preserving refinement is pinned by its ends, so transporting a cell's legs involves no
  choice.
\end{proof}

\section{Theorem C: at the point the polygraph is Artin's}\label{sec:thmC}

Since $Z$ is terminal, every serial wedge has one bipointed map to it and every wedge map commutes
over it, so $\Ch(Z) \cong \Wdg$: an object is a shape, a refinement a bare map of wedges. A run of
$Z$ is thus its event count; write $Z_N$ for the run on $N$ events. Both $\lorun e$ and $\hirun e$
are runs on $|e|$ events, so every cell of $Z$ is a loop at $Z_{|e|}$.

\begin{definition}\label{def:cuts}
  For a shape $a$ with $|a| = N$ put
  $\cut(a) := \{\, k \in [N-1] : \text{the events of strands } k, k+1 \text{ share a bead} \,\}$,
  and $\cut(d) := \cut(\text{shape of } d)$. For $i \neq k$ in $[N-1]$ let $\cox(i,k) := 3$ if
  $|i-k| = 1$ and $2$ otherwise, the Coxeter exponent of $S_N$.
\end{definition}

\begin{lemma}[{\cite[\S1.2]{EC1}}]\label{lem:cuts}
  $a \mapsto \cut(a)$ is a bijection from the shapes on $N$ events onto the subsets of $[N-1]$,
  with $|\cut(a)| = \deg a$. Hence $\Cell_n(Z_N,Z_N) \cong \{C \subseteq [N-1] : |C| = n\}$: the
  $1$-cells at $Z_N$ are indexed by $[N-1]$ --- write $\sigma_k$ for the one with $\cut = \{k\}$ ---
  and the $2$-cells by the pairs $i < k$. Moreover $|\Sq(e)| = \cox(i,k)$ for a degree-two $e$ with
  $\cut(e) = \{i<k\}$.
\end{lemma}
\begin{proof}
  A shape on $N$ events is a composition of $N$ and $\cut(a)$ is the complement of its bead
  boundaries in $[N-1]$, so the first claim is the classical correspondence and
  $|\cut(a)| = N - \#\text{beads} = \deg a$. For the last, $e$ has a bead of dimension $3$ exactly
  when three consecutive strands share one, that is when $k = i+1$; then $|\Sq(e)| = 3$ and
  otherwise $2$ (\Cref{prop:polygon}), matching $\cox$.
\end{proof}

\begin{lemma}\label{lem:pos}
  For $s \in \Steps(x)$ put
  $\pos_x(s) := \min\bigl(\step_x(\mathrm{lo}\,s), \step_x(\mathrm{hi}\,s)\bigr)$, so that $x$
  passes $s$ between times $\pos_x(s)$ and $\pos_x(s)+1$ (\Cref{lem:steps}). For any chain $d$: (i) $\pos_x : \Steps(x) \to \cut(d)$ is a bijection, for every
  $x \in \RunOver(d)$; (ii) $\pos_{x^s}(s) = \pos_x(s)$; (iii)
  $\pos_{\bot_d}(s) = \str(\mathrm{lo}\,s)$. If $d$ is a chain of $Z$ then moreover
  $\gamma(x \lessdot y) = \sigma_{\pos_x(s)}$ for a cover with step $s$.
\end{lemma}
\begin{proof}
  (i) The fibres of $\blk$ on the firing order are consecutive intervals of lengths
  $a_1,\dots,a_r$ (\Cref{lem:firing-blocks}), which do not depend on $x$; so the times $t$ with $t$
  and $t+1$ in one bead are $\cut(d)$ for every $x$. A step is a cover inside a bead
  (\Cref{lem:steps}), so $\pos_x$ lands there and is injective, distinct covers of a chain having
  distinct lower ends; and both sides have $\deg d$ elements (\Cref{lem:steps}, \Cref{lem:cuts}).
  (ii) Flipping exchanges the times of the two events and changes no other. (iii) $\bot_d$ fires in
  strand order. For the last claim, $\gamma(x \lessdot y) = m_{\{s\}}$, whose beads are the blocks
  of the firing order cut at every cover but $s$ (\Cref{prop:factorisations}): one bead of
  dimension $2$ at times $\pos_x(s), \pos_x(s)+1$ and the rest edges, so its cut set is
  $\{\pos_x(s)\}$.
\end{proof}

\begin{proposition}\label{prop:alternate}
  Let $e$ be a chain of $Z$ with $\deg e = 2$ and $\cut(e) = \{i<k\}$, and let
  $\bot_e = x_0 \lessdot \cdots \lessdot x_{|\Sq(e)|} = \top_e$ be the sweep of the standard order,
  with $s_t$ the $(t+1)$-st square. Then $\pos_{x_t}(s_t)$ is $i$ for $t$ even and $k$ for $t$ odd,
  so $\src(e) = \sigma_i\sigma_k\sigma_i\cdots$ and $\tgt(e) = \sigma_k\sigma_i\sigma_k\cdots$,
  both of length $\cox(i,k)$.
\end{proposition}
\begin{proof}
  By \Cref{lem:pos}(i) each $x_t$ has exactly two steps, passed at $i$ and at $k$. For $t = 0$:
  $\pos_{\bot_e}(s) = \str(\mathrm{lo}\,s)$ and the standard order sorts by $\str(\mathrm{lo}\,-)$
  first, so if $e$ is one bead of dimension $3$ on strands $i,i+1,i+2$ its squares in order are
  $(i,i{+}1), (i,i{+}2), (i{+}1,i{+}2)$, with lower strands $i,i,i+1=k$; and if $e$ is two beads of
  dimension $2$ on $\{i,i{+}1\}$ and $\{k,k{+}1\}$ they are $(i,i{+}1), (k,k{+}1)$, with lower
  strands $i,k$. Either way the least square is passed at $i$ and the greatest at $k$. For the
  step: $s_t$ and $s_{t+1}$ are both steps of $x_{t+1}$ --- the first because a flip of a step is a
  step, the second by \Cref{lem:orderpath} --- and $\pos_{x_{t+1}}(s_t) = \pos_{x_t}(s_t)$ by
  \Cref{lem:pos}(ii), so $\pos_{x_{t+1}}(s_{t+1})$ is whichever of $i,k$ that is not. The words
  follow by \Cref{lem:pos}, and their length is $|\Sq(e)| = \cox(i,k)$ (\Cref{lem:cuts}). The dual
  order reverses the standard one, so its least square is the standard order's greatest, passed at
  $k$, and the same induction applies.
\end{proof}

\begin{definition}\label{def:artin-poly}
  \textbf{Artin's polygraph} $\Pgraph_{\mathrm{Artin}}$ has one $0$-cell per $N \in \N$; at $N$,
  the loops $\sigma_1,\dots,\sigma_{N-1}$; and for each pair $i<k$ in $[N-1]$ one $2$-cell with the
  words $\sigma_i\sigma_k\sigma_i\cdots$ and $\sigma_k\sigma_i\sigma_k\cdots$ of length
  $\cox(i,k)$, that is $\sigma_i\sigma_{i+1}\sigma_i = \sigma_{i+1}\sigma_i\sigma_{i+1}$ when
  $k = i+1$ and $\sigma_i\sigma_k = \sigma_k\sigma_i$ otherwise. These are Artin's relations
  \cite{Artin}, so the words at $N$ modulo the $2$-cells are the positive braid monoid $B_N^+$;
  since path composition reverses concatenation, $\End(N) = (B_N^+)\op$ and
  $\pres{\Pgraph_{\mathrm{Artin}}} = \BraidPosOp$.
\end{definition}

\begin{theorem}[Theorem C]\label{thm:C}
  $\Pgraph(Z) \cong \Pgraph_{\mathrm{Artin}}$, cell for cell and word for word.
\end{theorem}
\begin{proof}
  $\Run(Z) \cong \N$, and every cell of $Z$ is a loop, so both polygraphs are disjoint unions over
  $N$ of loop polygraphs and it suffices to compare at each $N$. There $\Cell_1 \cong [N-1]$ and
  $\Cell_2 \cong \{i<k\}$ by \Cref{lem:cuts}, matching Artin's generators and relations, and by
  \Cref{prop:alternate} the two words of the $2$-cell at $\{i<k\}$ are Artin's two words for that
  pair. A morphism of polygraphs bijective in every dimension is an isomorphism
  (\Cref{prop:present-univ}).
\end{proof}

\begin{corollary}\label{cor:brZ}
  $\Br(Z) \simeq \BraidPosOp$, so $\Br(Z)$ is the positive braid category up to word reversal.
  Moreover the terminal map gives $\Br(!_K) : \Br(K) \to \Br(Z)$, a braid-valued invariant computed
  cell by cell: a degree-one chain of $K$ goes to $\sigma_k$, where $k$ is the position at which its
  run passes its one step, and a degree-two chain goes to the Artin relation of its two positions.
  This is the promised sharpening of \cite{PZ}: the braid is read off the chains of $K$ directly,
  with no nerve and no realization. Note that no braid theory was used to get it --- in particular
  Matsumoto's theorem, which identifies $B_N^+$ with the Garside germ, is not needed.
\end{corollary}
\begin{proof}
  \Cref{thm:A} at $K = Z$, \Cref{thm:C}, and \Cref{thm:B} at $!_K$.
\end{proof}

\printbibliography

\end{document}
